\section{Expansion in one length}
\label{ell:length-expansion}

Fix an edge length $L$ that is still finite.  Open that edge,
retaining all other lengths and ordinary parameters.  Let $Y_0$
denote the resulting solution, including compact unknowns and
any finite triangular system already appended.  At the two
new ends it is initially decaying.  We will close those ends
by adding finitely many terms of the form
\begin{equation}
                       e^{-pL}L^jY_{pj}.
 \label{ell:length-terms}
\end{equation}
The coefficient fields are solutions of the disconnected
triangular linearization.  They may grow at the two new ends.
If their largest growth exponent there is $b$, assign the term
the usable order
\begin{equation}
                 s=p-\max(0,b).
 \label{ell:usable-order}
\end{equation}
Take the maximum over both ends.  Polynomial powers in $\tau$
or $L$ do not change this order.  All $p,s$ are on the grid
$\lambda\Z$, because the tail powers and indicial roots are
integers and the seam prescription uses $\lambda=1/n$.
The terms constructed below have $s>0$.

Coefficient fields decay on every other still-finite half-end and
may have prescribed finite growth on older opened ends. At the
pair now being closed, the scalar factor in each length term
compensates its incoming growth through the usable order $s$.
The remainder estimate will keep these three bounds separate.

\begin{lemma}[Finite length expansion]
\label{ell:one-length}
For every desired exponential precision, the actual small
solution and every appended field have an expansion by finitely
many terms in Equation~\eqref{ell:length-terms}, with an
exponentially better remainder.  It is uniform in the remaining
finite lengths on their quadratic regions, in a fixed smaller
ordinary-parameter polydisc.  For older opened ends the conclusion
holds in every specified finite tail norm.  Coefficient fields
are obtained by the finite constructions in
Lemmas~\ref{ell:prescribed-tails} and~\ref{ell:triangular}.
\end{lemma}

\begin{proof}
We give the induction and then its remainder estimate.  Process
usable orders $s=\lambda,2\lambda,\ldots$.  Products add lower
bounds for usable order, since
\[
 \max(0,b_1+b_2)\le\max(0,b_1)+\max(0,b_2).
\]
No initial bulk residual is present: $Y_0$ solves the opened
equations.  Expand the remaining bulk equations by a finite
Taylor formula about $Y_0$, placing their full linearization
on the left.  Core equations, compact coefficients, and
previously imposed asymptotic conditions are part of these
equations.  At order $s$ only finitely many products of the
earlier positive-order terms occur.  For each coefficient of
$e^{-pL}$ times a polynomial in $L$, cancel its complete forced
field by the disconnected linearized inverse, with zero free
incoming constants on the two new ends.  A source assigned
order $s$ has growth at most $p-s$ there; by
Lemmas~\ref{ell:prescribed-tails} and~\ref{ell:triangular} its
correction has the same upper growth bound.  When this bound
is negative the inverse may also contribute decaying
homogeneous terms; these still have positive usable order.
More uniformly, record each individual product with its
product lower bound and allow the bound $\max(0,p-s)$.
For the terms at issue $p\ge s$, by
Equation~\eqref{ell:usable-order}.  The newly solved linear
terms vanish exactly.  Their later bulk interactions are
higher-order products.

Next evaluate the value, normal trace, and critical-test
residuals at $\tau=L$, using integer tails to finite precision.
At order $s$ they are finite harmonic sums times
$e^{-sL}$ and polynomials in $L$; include the prescribed
$e^{-\lambda L}z_e$ term.  For every incoming root $b\ge0$
there is a homogeneous disconnected variation with prescribed
leading mode $e^{b\tau}$ on its selected side.  Prescribe zero
free nondecaying constants elsewhere.  Its difference from
that leading mode has strictly smaller growth at the two
new ends, because the coefficients coupling channels decay
there.  For $b=0$ all other new-end components decay.
Multiply this variation by $e^{-(s+b)L}$ times a polynomial
in $L$.  Its leading seam data have order $s$.  The block
matrices in Lemma~\ref{ell:seam-inverse} solve every leading
seam residual.  Their lower tails contribute only later
orders.  Matching at all other finite seams is preserved by
the disconnected inverse.  This completes the induction step.
Finiteness follows from the positive minimum tag $\lambda$,
the finite harmonic support at each previously reached order,
and the integer root lists.  In particular, no angular
analyticity estimate uniform in all expansion orders is used.

We now justify that the construction approximates the actual
fields.  Restrict each coefficient to the new halves
$\Re\tau\le\Re L$.  Before applying the basic weight, a term
with usable order $s$ has field norm, or source norm with two
fewer derivatives, bounded by
\begin{equation}
            C(1+\Re L)^C e^{-s\Re L}.
 \label{ell:tag-bound}
\end{equation}
For a decaying coefficient this uses its whole half-cylinder
norm; for a growing coefficient it uses the finite tail and
$e^{b\Re\tau}\le e^{b\Re L}$.  Products consequently cost
only polynomials in this unweighted estimate.  Analytic Taylor
remainders beyond a fixed tag cutoff contain arbitrarily many
positive-order factors if the Taylor degree is chosen large
enough.  The base nonlinear field and its corrections on these
halves are small, so those Taylor estimates are valid.
Apply the basic weight only after the product estimate; it
costs $e^{\nu\Re L}$ once, not once per factor.  The fixed
polynomial source strengthening costs an arbitrarily small
additional exponential rate.  The same reasoning applies
to seam tails and their first normal traces, which have been
expanded to a strictly higher rate before evaluation.
It follows that a sufficiently high finite tag cutoff gives
any requested residual order in all basic norms.

On another still-finite half, the coefficient fields decay
in the fixed basic weight.  Residual bulk products contain
the two factors required by Lemma~\ref{ell:triangular},
which absorb the polynomial loss in that other length.
There is no seam residual there.  On an older opened end,
every term retains its full scalar factor $e^{-pL}$.
Products there are estimated to a higher finite tail precision
to offset the finite growth of the other factors.  For
unexpanded analytic operations on the base metric, its own
corrections at old opened ends remain decaying: their equations
and old-end conditions depend only on the base metric and
its own corrections.  Thus Taylor remainders there are
controlled by extracting additional finite Taylor powers;
growing auxiliary fields appear only in finite differential
products and triangular equations.  This proves residual
estimates in all required old-end tail norms as well.

Compare the nonlinear approximate metric first in the basic
decaying spaces, allowing its seam jumps as data.  The
small-branch inverse and the small nonlinear Lipschitz bound
give the same arbitrary exponential accuracy for its difference
from the true metric.  On an old opened end the difference
equation has flat principal operator, decaying coefficients
with integer tails, and the residual just bounded to higher
precision.  Apply Lemma~\ref{ell:flat-tail} repeatedly.  It
upgrades the initial small error to the same exponential
accuracy for any specified tail coefficients and remainder.
For an appended field, work in triangular order: its difference
equation has the same diagonal inverse and errors from previous
members multiplied by at most finite growth.  Apply
Lemma~\ref{ell:prescribed-tails}, also to differences of its
free prescriptions.  This proves the asserted comparison in
finite tail norms, without ever trying to put a growing
field itself in a decaying space.

All estimates use a fixed high Sobolev order.  A Taylor source
uses at most two derivatives and its inverse regains those
two derivatives.  Thus iteration creates no accumulating
differential loss.  Traces and contour shape differentiation
may use one lower order, and parameter derivatives use Cauchy
estimates in the independent analytic variables.
\end{proof}

